% Experiments. All numbers come from numbers.tex / numbers_results.tex and tables/*.tex.

\section{Experiments}
\label{sec:experiments}

\subsection{Setup}
\label{sec:setup}

We use UCF-Crime~\cite{sultani2018real} (290 test videos) and XD-Violence~\cite{wu2020xdviolence} (800 test videos) with the CLIP features, splits, and frame-level ground truth released with VadCLIP. Following VadCLIP, we report frame-level AUC of $s^{(1)}$ on UCF-Crime and frame-level AP of $p$ on XD-Violence. Controlled experiments and ablations use a test-independent fixed-epoch protocol with matched seeds: every model is trained for a fixed number of epochs and the final checkpoint is evaluated. Each condition is run with seeds 234, 235, and 236, which share initialization and sample order across arms, and we report mean$\pm$SD, the paired difference per seed, and the number of positive seeds. We additionally evaluate the final method using VadCLIP's official training script.

\subsection{Main Results}
\label{sec:main}

% Auto-generated by scripts/make_tables.py. Do not edit by hand.
\begin{table}[t]
\centering
\small
\setlength{\tabcolsep}{4pt}
\caption{Results on UCF-Crime (AUC, \%) and XD-Violence (AP, \%) with CLIP features. Published rows are taken from VadCLIP~\cite{wu2024vadclip} and serve as benchmark context; each follows its own training protocol. Method claims rest on the matched comparison: mean$\pm$SD over three seeds, paired $\Delta$ per seed, number of positive seeds in parentheses. The bottom block repeats the matched comparison under VadCLIP's official script, which selects the checkpoint on the test set; per-seed values are in Table~\ref{tab:official}. Only rows within a block are directly comparable.}
\label{tab:main}
\begin{tabular}{lcc}
\toprule
Method & UCF AUC & XD AP \\
\midrule
RTFM~\cite{tian2021rtfm} & 85.66 & 78.27 \\
UR-DMU~\cite{zhou2023urdmu} & 86.75 & 82.41 \\
CLIP-TSA~\cite{joo2023cliptsa} & 87.58 & 82.17 \\
VadCLIP~\cite{wu2024vadclip} & 88.02 & 84.51 \\
\midrule
\multicolumn{3}{l}{\emph{Matched comparison (same protocol and seeds)}} \\
VadCLIP, retrained & 86.25$\pm$0.83 & 77.57$\pm$2.19 \\
Ours & 86.53$\pm$0.85 & 79.23$\pm$1.87 \\
Paired $\Delta$ & +0.28$\pm$0.15 (3/3) & +1.66$\pm$0.34 (3/3) \\
\midrule
\multicolumn{3}{l}{\emph{VadCLIP's official script (best-epoch selection)}} \\
VadCLIP, retrained & -- & 82.91$\pm$1.23 \\
Ours & -- & 83.63$\pm$1.18 \\
\bottomrule
\end{tabular}
\end{table}


Table~\ref{tab:main} places our results in the context of published CLIP-feature methods and reports the matched comparison. Against VadCLIP retrained under the same protocol and seeds, our method improves UCF-Crime AUC by \ucfGainMean{}$\pm$\ucfGainSd{} points and XD-Violence AP by \xdGainMean{}$\pm$\xdGainSd{} points, and the paired gain is positive on every seed for both datasets. The two gains come from different designs: Stage~2 leaves $s^{(1)}$ unchanged, so the UCF-Crime gain is produced by Stage~1, whereas Stage~1 alone changes XD-Violence AP by only \xdStageOneMean{} points and the XD-Violence gain is produced by Stage~2.

\subsection{Mechanism Ablations}
\label{sec:ablation}

% Auto-generated by scripts/make_tables.py. Do not edit by hand.
\begin{table}[t]
\centering
\small
\setlength{\tabcolsep}{4pt}
\caption{Mechanism ablations (\%, mean$\pm$SD over three matched seeds). $\Delta$ is the paired difference to VadCLIP (B0) unless stated otherwise. (a) Post-aggregation computes the difference on the pooled sequence with all other settings fixed. (b) C runs the same budget from the same checkpoint with $\lambda=0$.}
\label{tab:ablation}
\begin{tabular}{lcc}
\toprule
Variant & Score & $\Delta$ \\
\midrule
\multicolumn{3}{l}{\emph{(a) Stage 1, UCF-Crime AUC}} \\
VadCLIP (B0) & 86.25$\pm$0.83 & -- \\
Post-aggregation variation & 85.86$\pm$0.84 & $-$0.39$\pm$0.10 (0/3) \\
Pre-aggregation variation & 86.53$\pm$0.85 & +0.28$\pm$0.15 (3/3) \\
\quad Pre $-$ Post &  & +0.67$\pm$0.05 (3/3) \\
\midrule
\multicolumn{3}{l}{\emph{(b) Stage 2, XD-Violence AP}} \\
VadCLIP (B0) & 77.57$\pm$2.19 & -- \\
Stage 1 & 77.44$\pm$2.52 & $-$0.14$\pm$0.33 (1/3) \\
+ CLASM only (C) & 77.59$\pm$2.28 & +0.02$\pm$0.30 (2/3) \\
+ Restricted objective & 79.23$\pm$1.87 & +1.66$\pm$0.34 (3/3) \\
\quad Restricted $-$ C &  & +1.64$\pm$0.43 (3/3) \\
B0 + restricted objective & 79.46$\pm$1.48 & +1.89$\pm$0.76 (3/3) \\
\bottomrule
\end{tabular}
\end{table}


\textbf{Where temporal variation is extracted.} Table~\ref{tab:ablation}(a) computes the same difference after aggregation, $D'_j = \lvert A_j - A_{j-1} \rvert$, with every other setting fixed. This comparison isolates the position of the difference relative to pooling. Post-aggregation variation falls below the baseline on every seed (\postBaseMean{} points), while pre-aggregation variation exceeds it by \prePostMean{}$\pm$\prePostSd{} AUC points (\prePostPos{} seeds). The benefit therefore comes from preserving within-segment change, not from adding a difference stream as such.

\textbf{Where semantic optimization is applied.} Table~\ref{tab:ablation}(b) compares Stage~2 with an equal-budget control C that runs the same epochs from the same checkpoint and sample order with $\lambda=0$. The restricted semantic objective adds \tMinusCMean{}$\pm$\tMinusCSd{} AP points over C (\tMinusCPos{} seeds), so the gain comes from the objective rather than from extra training. Applied directly to the retrained VadCLIP, Stage~2 improves AP by \bzTMean{}$\pm$\bzTSd{} points (\bzTPos{} seeds), which attributes the XD-Violence gain to Stage~2; placing Stage~1 before it changes AP by a further \tMinusBzTMean{} points. During Stage~2, $s^{(1)}$ on a fixed probe batch did not change and no parameter outside $\mathrm{mlp}_1$ was updated.

\subsection{Robustness and Efficiency}
\label{sec:robust}

Under VadCLIP's official training setting, XD-Violence AP improves from \offBase{}$\pm$\offBaseSd{} to \offOurs{}$\pm$\offOursSd{}, with a paired gain of \offGainMean{}$\pm$\offGainSd{} points and \offGainPos{} positive runs. The improvement therefore persists beyond our protocol. Absolute values differ between the two settings, so each comparison is made within one setting.

The method adds \pvrParams{} parameters, increases inference latency by \cxUcfLatRel--\cxXdLatRel\%, and adds at most \cxUcfMemDelta{}~MB peak memory, measured on an \cxGpu{} per window of 256 snippets. This overhead keeps the detector suitable for lightweight surveillance deployment.
